\section{民事法律行为}
\begin{enumerate}
  \item 民事法律事实。
        \begin{enumerate}
          \item 事件。与人的意志无关。\al{洪水、海啸、人的出生}。
          \item 行为。
                \begin{enumerate}
                  \item 当事人决定。能够发生，意欲发生的，主动追求的法律后果 \yjt \zd{民事法律行为}。\al{合同、婚姻、遗嘱}。
                  \item 法律决定。\zd{事实行为}，不考虑行为人的主观意图。\al{侵权、无因管理、创作}。
                \end{enumerate}
        \end{enumerate}
  \item 法定分类。
        \begin{enumerate}
          \item 双方民事法律行为。\al{合同}。
          \item 单方民事法律行为。\al{动产所有权的抛弃}。
        \end{enumerate}
  \item 学理分类。
        \begin{enumerate}
          \item 财产行为。\al{遗嘱}。
                \begin{enumerate}
                  \item \zd{负担行为}。只发生债权债务关系。\al{把房子租给你}。
                  \item \zd{处分行为}。所有权变动。\al{物权行为}。
                \end{enumerate}
          \item 身份行为。\al{婚姻、收养}。
        \end{enumerate}
  \item 成立和生效。成立 \yjt 效力 \yjt 生效 \yjt 失效。\zd{严格区分}。
        \begin{enumerate}
          \item \al{甲卖表给乙，乙不买}。
          \item \al{甲卖毒品给乙，乙买}。
          \item \al{考试第 1 名，给你买玩具}。
          \item \al{超市买瓶矿泉水}。
        \end{enumerate}
  \item 效力。\zd{432 结构}。
        \begin{enumerate}
          \item 合同。效力最完整。有效；效力瑕疵 \yjt 无效、效力待定、可撤销。
                \begin{enumerate}
                  \item 有效。\zd{强、公、主}体具有相应的民事行为能力、\zd{意}。
                  \item 无效。\zd{强公主意恶}。\al{20 岁甲和乙结婚；代孕}。
                        \begin{enumerate}
                          \item \zd{双方通谋}虚假意思。阳合同/条款无效（\zd{识别哪个是假的}），阴合同具体分析。
                          \item \zd{恶意串通}。互相勾结、谋取私利、损害他人合法权益。结合代理考。\al{私下，回扣}。
                          \item 法律后果。不是不发生任何法律后果。\al{甲被乙打成重伤，达成协议支付医药费 5W，不得告发}。
                        \end{enumerate}
                  \item 效力待定。\zd{只能发生在合同里}。
                        \begin{enumerate}
                          \item 民事行为能力欠缺。\al{未成年打赏网络主播}。
                          \item 处分能力欠缺。\al{无权代理}。
                        \end{enumerate}
                  \item 可撤销。\zd{本质是意思表示不自由、不真实}，赋予特定权利人来最终决定。考为什么可撤销。
                        \begin{enumerate}
                          \item 重大误解。\zd{自身原因 + 主给付义务 + 认识错误}，双方均有撤销权。
                                \begin{itemize}
                                  \item 重大。\zd{主给付义务}的认识错误。\al{付 5 块输成 500}。
                                  \item \zd{从给付义务}。辅助合同目的的实现，\zd{违反不构成重大误解}。\zd{配套设施、单证资料}。
                                  \item 附随义务（诚信原则）。辅助主、从给付义务实现。\al{通知、协助、保密}。
                                  \item 误解。行为的性质，\zd{勿把有偿当无偿}，\al{酒店房间里的 XO}
                                  \item 对方当事人。不包括对行为能力的认识。\al{未成年}。
                                  \item 标的物。品种、规格、数量、价格、\zd{第三人传达错误}。\al{网店标错价格，薅羊毛}；\al{古玩、字画风险自担}。
                                \end{itemize}
                          \item 欺诈（同刑法欺诈）。\zd{外部原因 + 主/从给付义务 + 认识错误}。
                                \begin{itemize}
                                  \item \al{人造钻石冒充天然钻石；凶宅就不告诉你；知假卖假}。
                                  \item \zd{相对方是否知情}。喜欢考第三人欺诈。\al{托、双簧}。
                                \end{itemize}
                          \item 胁迫（同刑法敲诈勒索）。手段 \hz \zd{目的}不正当，非法性、不法性、不正当性，谁实施在所不问。
                                \al{不让退班起诉你；甲盗得 1W，跑路途中神色慌张，乙暴力胁迫不给钱就报案}。
                          \item 显失公平。\zd{主客观相结合}。相对的重点。
                                \begin{itemize}
                                  \item 主观。\zd{利用}危困状态、缺乏判断能力等。\al{甲喝醉酒打车，司机说 80，甲输错付了 8000}
                                  \item 客观。造成双方权利义务不对等，显著失衡。
                                  \item 单方允诺。\al{假一罚十}。
                                \end{itemize}
                          \item \zd{未撤销前有效}。或诉或裁，双重除斥期间限制（主观起算）。
                                \begin{itemize}
                                  \item 重大误解 90 天，其余 1 年。
                                  \item 主观起算。胁迫自胁迫行为终止之日起（公平原则）。
                                  \item \zd{最长保护期}。自民事法律行为发生之日起 \zd{5} 年内。
                                  \item \zd{主客观同时起算，到期在先原则}。
                                \end{itemize}
                        \end{enumerate}
                \end{enumerate}
          \item 婚姻。有效、无效、可撤销。
          \item 遗嘱。有效、无效。
        \end{enumerate}
  \item 特殊的民事法律行为。\zd{影响的是生效、失效这两个阶段}。常考点。
        \begin{enumerate}
          \item 附条件。条件（这事儿）\zd{发生与否不确定}。
                \begin{enumerate}
                  \item \zd{不可能条件}。具体情况具体分析 \yi 无效；\er 视为未附条件。\al{太阳西边出来送你手表；不续租你房}。
                  \item 不能不正当的促进/阻止条件成就，否则\zd{视为条件不/已成就}（拟制效力）。本质是欺诈行为。
                \end{enumerate}
          \item 附期限。\zd{期限的到来是必然的}。考过 1 次。
        \end{enumerate}
\end{enumerate}
